\section{Problem Formulation}
\label{sec:problem}

\subsection{System and Trust Model}
\label{subsec:trust-model}

We consider a blockchain-based healthcare data-sharing system with two user roles, following the consent model of Jaiman and Urovi~\cite{jaiman2020consent}: \textit{data providers}, who provide protected resources and govern access through consent policies, and \textit{data requesters}, who seek access to those resources subject to the applicable policies.
TPS-ZKPI mediates this interaction through three components: a TPS operating off-chain, an on-chain access-control manager (ACM), and a data access layer.
The provider, which may be the data subject or an organization holding the data~\cite{jaiman2020consent}, defines a consent policy $P$ for a protected resource and registers an on-chain commitment $h$ to the encoded policy.
The requester seeks access by submitting policy-relevant authorization evidence, such as attributes, credentials, and access context.
The TPS constructs the private witness and generates a zero-knowledge proof
of policy compliance.
The ACM verifies the proof on-chain and records the authorization outcome, while the data access layer grants access according to the recorded decision.

TPS-ZKPI relies on the following assumptions.
First, the provider defines the consent policy and commits to the corresponding policy instance on-chain.
Second, the TPS is trusted to preserve the confidentiality of policy-bound proving material and requester-submitted authorization evidence, and to follow the prescribed proof-generation workflow.
Third, the commitment scheme is assumed to be hiding and binding, and the zero-knowledge proof system is assumed to satisfy completeness, soundness, and zero knowledge with respect to the implemented policy-compliance relation.
Issuer governance, credential issuance and revocation, and real-world identity verification are treated as external lifecycle functions.

Conceptually, the TPS is an independent proof service, distinct from the provider; this separation is what makes relocating trust meaningful, since it prevents the provider from being the single party that holds both the policy and the requester evidence. Possible instantiations of such an independent service include a third-party proof service, a TEE-backed proof service that makes the trust assumption hardware-attested, or MPC/collaborative SNARK proof generation in which no single party sees both policy material and requester evidence. For evaluation simplicity, the current evaluation co-locates the TPS with the provider (a provider-operated TPS).
This co-location concentrates trust in the provider, which holds both the policy and the requester's authorization evidence.
A compromised provider-operated TPS could substitute inputs if their origin is not independently authenticated, or refuse to generate proofs for legitimate requesters.

The policy commitment and on-chain proof verification ensure that accepted proofs demonstrate satisfaction of the committed policy relation, regardless of how the TPS is instantiated.
However, they do not guarantee the provenance of evidence that is not issuer-signed, which remains dependent on the TPS (Table~\ref{tab:threat-guarantee}).
Policy confidentiality, in contrast, is not cryptographically enforced against the TPS itself; it rests on this trust assumption.

Table~\ref{tab:threat-guarantee} summarizes, for each actor or failure mode, how each of four security properties (policy confidentiality, authorization correctness, liveness, and decision auditability) holds: whether it is enforced cryptographically, supported only empirically, dependent on the TPS trust assumption, left to external functions, or not provided by the prototype.
% Threat/guarantee matrix for the trust model of Section 3.1 (reviewer request).
% Rows: parties or failure modes. Columns: the four guarantees the workflow provides.
% Cell values: C = enforced cryptographically (commitment binding, Groth16 soundness/ZK,
%   EdDSA-Poseidon unforgeability, on-chain checks); T = holds only if the TPS follows
%   the prescribed workflow; E = external/lifecycle function; -- = not provided.
\begin{table}[htbp]
  \centering\footnotesize
  \caption{What each guarantee of TPS-ZKPI rests on, per threat. C: cryptographic
    (commitment binding, proof soundness and zero knowledge, credential signature,
    on-chain checks); T: trust in the TPS following the prescribed workflow; E: external
    or lifecycle function; --: not provided by the prototype.}
  \label{tab:threat-guarantee}
  \setlength{\tabcolsep}{4pt}
  \begin{tabular}{p{0.18\linewidth}p{0.20\linewidth}p{0.22\linewidth}p{0.13\linewidth}p{0.16\linewidth}}
    \toprule
    Threat & Policy confidentiality & Authorization correctness & Liveness & Accountability \\
    \midrule
    Passive public-chain observer & C (proof, commitments) + empirical (gas padding) & C & n/a & C (recorded outcome) \\
    Requester (untrusted) & C: sees only the public view & C: cannot forge a proof or reuse one in another context & n/a & C \\
    Honest-but-curious TPS & T: TPS holds plaintext policy and evidence & C & T & -- \\
    Compromised or colluding TPS & --: policy and evidence exposed & C for committed relation; T for origin of unsigned evidence & --: may refuse to prove & --: refusal indistinguishable from denial \\
    Provider-operated TPS (this evaluation) & T, concentrated at the provider & C & T & -- \\
    Forged or unauthorized issuer & n/a & C against forgery under the committed key; E for issuer legitimacy and revocation & n/a & E \\
    \bottomrule
  \end{tabular}
\end{table}


Attribute-origin guarantees depend on the evidence type.
For DUO terms evaluated from the claims in an issuer-signed profile credential presented by the requester (hereafter, profile-backed terms), the prototype verifies the credential's signature under an issuer key selected by the provider and bound to the policy commitment.
The signature scheme, EdDSA-Poseidon, is the Edwards-curve Digital Signature Algorithm (EdDSA)~\cite{bernstein2012high} over the Baby Jubjub curve~\cite{whitehat2020babyjubjub}, instantiated with Poseidon~\cite{grassi2021poseidon}, a hash function designed for efficient evaluation inside zero-knowledge circuits.
The credential binds the policy commitment, requester pseudonym, protected resource, claims, validity interval, and credential identifier; the corresponding policy attributes are derived inside the circuit.
Thus, subject to signature unforgeability and Poseidon collision resistance, these values cannot be modified without invalidating the proof.

\subsection{Public-Chain Adversary}

The primary adversary is a passive public-chain observer.
The adversary can inspect smart-contract code and monitor public blockchain artifacts across multiple authorization sessions, including policy commitments, submitted proofs, public inputs, emitted events, calldata, timing, and gas usage.
The adversary does not observe provider-to-TPS policy material, commitment-opening randomness, requester-submitted authorization evidence, credential witnesses, or the TPS's witness construction process.
These data and computations remain off-chain and are not exposed in the public-chain transcript.
We assume the off-chain channels themselves are confidential (e.g., TLS-protected).

Each policy-enforcement instance corresponds to an authorization session, in
which a requester seeks access to a protected resource under a committed policy.
For each such session, the adversary observes a public view
\begin{equation}
	\mathit{view}_{\mathrm{pub}}= (h, \pi, o, \mathit{pp}, d_{\chi}, \tau, \mathit{id}_{\mathrm{req}}, \mathit{id}_{\mathrm{res}}, \nu, \mathit{meta}),
\end{equation}
where $h$ is the on-chain commitment to the hidden policy, $\pi$ is the zero-knowledge proof submitted for on-chain verification, $o$, $\mathit{pp}$, $d_{\chi}$, and $\tau$ are the remaining public signals of the proof statement (the authorization output, a parameter commitment, a context digest, and a request epoch; Section~\ref{sec:design}), $\mathit{id}_{\mathrm{req}}$, $\mathit{id}_{\mathrm{res}}$, and $\nu$ are the requester, resource, and session-nonce values carried in the access transaction, and $\mathit{meta}$ denotes transaction-level metadata
exposed by the enforcement layer, such as calldata footprint, event structure, public-input layout, timing, and gas usage. The adversary's goal is to infer hidden information about the committed policy or the authorization interaction from this public view.

In addition to the public-chain observer, we consider the requester as untrusted with respect to policy confidentiality.
The requester initiates access requests, observes its own submitted evidence and authorization outcome, and inspects the public-chain transcript.
However, the requester does not receive the encoded policy, commitment-opening randomness, or policy-bound proving material, which remain within the provider-to-TPS proof-generation boundary.
This restriction matters because a requester that learned the policy could infer sensitive facts about the data subject from the consent terms themselves and, in addition, would learn the policy's full acceptance region (which attribute combinations are granted), enabling targeted evidence construction or probing beyond its own legitimate sessions.

\subsection{Privacy Objectives}
\label{subsec:privacy-objectives}

We operationalize consent policy privacy as three measurable, attack-oriented objectives that cover the information a public-chain observer could extract about a hidden policy: its semantic content, which of several similar policies it is, and its structural form.
These correspond to Policy Content Hiding, Policy Indistinguishability, and Structure and Parameter Privacy, respectively.

\textit{Policy Content Hiding (PCH).}
PCH requires that public-chain artifacts should not reveal the semantic content of the hidden consent policy, such as data-use terms, disease categories, research purposes, institutional restrictions, or ethics-related requirements.

\textit{Policy Indistinguishability (PI).}
PI requires that different candidate policies should not become distinguishable from their public verification artifacts when they are compiled into the same fixed public interface.
This captures whether two policies can be told apart from commitments, proofs, public
parameters, or metadata even when their direct representations are hidden, that is,
whether an observer can tell which policy governs an execution.

\textit{Structure and Parameter Privacy (SPP).}
SPP requires that public-chain artifacts should not reveal structural or parametric properties of the
hidden policy, such as logical structure, nesting depth, predicate count, or set size.

\subsection{Problem Statement}

Given $P$, requester authorization evidence, a globally observable public blockchain, and the TPS trust assumption stated in Section~\ref{subsec:trust-model}, the problem is to enforce access to a protected healthcare resource subject to four requirements.
First, authorization must be publicly verifiable.
The ACM and any observer must be able to verify that the requester satisfies the committed policy without learning the policy or executing it in plaintext.
An accepted proof must be bound to the committed policy and the request's authorization context, including the requester pseudonym, resource, and session, to prevent reuse in a different context.
Second, the policy must remain confidential.
Public-chain artifacts should protect policy content, structure, and parameters and preserve policy indistinguishability as defined by PCH, SPP, and PI (Section~\ref{subsec:privacy-objectives}).
Requesters should learn no policy information beyond what their own authorization outcomes reveal, and requester evidence should remain off-chain.
Third, authorization decisions should be auditable.
Each processed request, whether granted or denied, should be recorded on-chain together with its bound authorization context, allowing the decision to be verified later without revealing the policy logic.
Fourth, one policy-independent verification interface should serve every policy in the supported class, so that policies can be added or changed without changing the on-chain verifier.
